\section{Limiti e colimiti nella pratica matematica}
perché si chiamano limiti?
\begin{remark}
 \Todo{Limiti in analisi come limiti categoriali}
\end{remark}
\subsection{Limiti in algebra, topologia e logica}
\begin{example}\label{ex_lim_kernel}
 Il \emph{nucleo} di una applicazione lineare \(f \colon V \to W\) tra spazi vettoriali (o tra gruppi abeliani, o \(R\)-moduli,\dots), già introdotto in \ref{ker_e_coker}, è un esempio di limite: si tratta dell'equalizzatore della coppia
 % https://q.uiver.app/#q=WzAsMixbMCwwLCJWIl0sWzEsMCwiVyJdLFswLDEsIjAiLDAseyJvZmZzZXQiOi0yfV0sWzAsMSwiZiIsMix7Im9mZnNldCI6MX1dXQ==
 \[\begin{tikzcd}[cramped]
   V & W
   \arrow["0", shift left=2, from=1-1, to=1-2]
   \arrow["f"', shift right, from=1-1, to=1-2]
  \end{tikzcd}\]
 In effetti una nozione di `nucleo' in questo senso si può dare in ogni categoria che abbia un \emph{oggetto zero}, nel senso di un oggetto che sia allo stesso tempo iniziale e terminale (questo è il caso di ciascuna delle categorie appena menzionate, ma anche della categoria degli insiemi puntati di \ref{}). In tal caso, si ricordi che la `mappa zero' consiste della composizione \(0_{VW} \colon V \to 0 \to W\) e in un certo senso evidente \(\ker f\) è in quel caso il massimo sottoggetto \(i_{\ker f} \colon \ker f \mono V\) tale che \(f\circ i_{\ker f} = 0_{VW}\). Nel caso degli insiemi puntati, dove \(V=(A,a_0)\), si tratta del sottoinsieme \(K\subseteq A\) determinato come \(K=\{a\in A\mid fa = b_0\}\), se \(W=(B,b_0)\) è l'insieme puntato codominio.
\end{example}
\begin{example}\label{exa_sottoggetti_limite}\index{Funtore!--- dei sottoggetti}\index{Sottoggetto}
 Il funtore dei sottoggetti introdotto in \ref{def_fun_sub_onset} fa uso della nozione di prodotto fibrato in \(\ctSet\): infatti, data una funzione \(u \colon A\to B\) tra insiemi, la funzione \(\Sub(u) \colon \Sub(B) \to \Sub(A)\) si determina dal prodotto fibrato
 \[\begin{tikzcd}
   u^*E =\pull AEB \arrow[r] \arrow[d, hook] & E \arrow[d, hook] \\
   A \arrow[r, "u"'] & B
  \end{tikzcd}\]
 Infatti, la funzione \(\Sub(u)\) manda un sottoggetto \(E\mono B\) nel sottoggetto \(u^*E\mono A\).
\end{example}

\begin{example}\label{exa_grafico_limite}\index{Grafico!--- di una funzione}
 \Todo{Grafici di funzioni}
\end{example}

\begin{example}\label{exa_tabulatore_relazione_limite}\index{Tabulatore}\index{Relazione}
 \Todo{Tabulatore di una relazione}
\end{example}

\begin{example}\label{exa_punti_fissi_limite}\index{Punto fisso}\index{Limite!Punti fissi come ---e}
 \Todo{Punti fissi come limiti}
\end{example}

\begin{example}\label{exa_categoria_elementi_limite}\index{Categoria!--- degli elementi}
 \Todo{La categoria degli elementi come limite}
\end{example}

\begin{example}\label{exa_comma_limite}\index{Categoria!--- comma}
 \Todo{\((F/G)\) come limite}
\end{example}

\begin{example}\label{exa_nat_fg_limite}\index{Trasformazione naturale}
 \Todo{\(Nat(F,G)\) come limite}
\end{example}

\begin{example}\label{exa_spazi_lacci_limite}\index{Spazio!--- di lacci}
 \Todo{Spazi di lacci come limite}
\end{example}

\begin{example}\label{exa_costruzione_limiti_cat}
 \Todo{Costruzione dei limiti in \(\ctCat\)}
\end{example}

\begin{example}\label{exa_estensione_destra_lungo_g}\index{Estensione!--- destra}
 \Todo{Estensione destra di F lungo G}
\end{example}

\begin{examples}\label{exas_estensioni_destre}
 Esempi di estensioni destre
\end{examples}
\subsection{Colimiti in algebra, topologia e logica}
\begin{example}\label{exa_prodotto_semidiretto_gruppi_monoidi}\index{Prodotto semidiretto!--- di gruppi e monoidi}\index{Gruppo!--- prodotto semidiretto}\index{Monoide!--- prodotto semidiretto}
 \Todo{Prodotto semidiretto di gruppi e monoidi}
\end{example}

\begin{example}\label{exa_colimite_da_completare}\index{Colimite}
 \Todo{}
\end{example}

\begin{example}\label{exa_conuclei_quozienti}\index{Conucleo}\index{Quoziente!---i e conuclei}
 \Todo{Conuclei e quozienti}
\end{example}

\begin{example}\label{exa_spazi_orbite_colimiti}\index{Spazio!--- di orbite}\index{Colimite!spazio di orbite come ---}
 \Todo{Spazi di orbite come colimiti}
\end{example}

\begin{example}\label{exa_pi_zero_colimite}\index{\(\pi_0\)}\index{Colimite!\(\pi_0\) come ---}
 \Todo{\(\pi_0\) come colimite.}
\end{example}

\begin{example}\label{exa_categoria_elementi_colimite}\index{Categoria!--- degli elementi}
 \Todo{La categoria degli elementi come colimite}
\end{example}

\begin{example}\label{exa_c_star_d_colimite}\index{Giunto}
 \Todo{\(\ctC\star\ctD\) come colimite}
\end{example}

\begin{example}\label{exa_sospensione_colimite}\index{Sospensione!--- ridotta}\index{Colimite!--- sospensione ridotta come}
 \Todo{Sospensione e Sospensione ridotta come colimite}
\end{example}

\begin{example}\label{exa_costruzione_colimiti_cat}\index{Categoria!--- dei colimiti in \(\ctCat\)}\index{Colimite!--- costruzione in \(\ctCat\)}
 \Todo{Costruzione dei colimiti in \(\ctCat\)}
\end{example}

\begin{example}\label{exa_estensione_sinistra_lungo_g}\index{Estensione!--- sinistra}
 \Todo{Estensione sinistra di F lungo G}
\end{example}

\begin{examples}\label{exas_estensioni_sinistre}
 Esempi di estensioni sinistre
\end{examples}
\begin{theorem}
 Teorema di Knaster-Tarski (l'esistenza di un punto fisso è determinata da una `prop univ')
\end{theorem}
\subsection{Costruzioni più sofisticate che usano limiti}
\Todo{molti esempi disseminati per il capitolo 1 sono stati architettati per evitare di usare il linguaggio dei limiti;	ora diventano molto più snelli. Metto qui un po' di esempi senza pretesa di ordine (andranno spostati qua e là in ogni caso).}
\begin{example}\label{ex_N_coequalizzatore_ctCat}\index{Coequalizzatore!in \(\ctCat\)}\index{Gruppo libero!come coequalizzatore}
 \(\bbN\) come coequalizzatore in \(\ctCat\); più in generale, gruppo libero come coequalizzatore in \(\ctCat\).
\end{example}

\begin{example}\label{ex_coequalizzatori_riflessivi_ctMod}\index{Coequalizzatore!riflessivo}\index{\(\ctMod(\Omega,s)\)!coequalizzatori riflessivi}
 Coequalizzatori riflessivi in categorie della forma \(\ctMod(\Omega,s)\) costruiscono i coprodotti (e.g., in monoidi, gruppi, anelli, \dots).
\end{example}

\begin{proposition}\label{prop_categorie_filtrate_setacciate}\index{Categoria!filtrata}\index{Categoria!setacciata}
 \Todo{Categorie filtrate e setacciate}
\end{proposition}

\begin{definition}\label{def_funtore_sottooggetti}\index{Funtore!--- dei sottooggetti}\index{Sottoggetto!funtore dei ---}
 \Todo{Il funtore dei sottoggetti, stavolta per davvero}
\end{definition}

\begin{example}\label{ex_spiga_prefascio}\index{Spiga}\index{Fascio!spiga}\index{Prefascio!spiga}
 Spiga di un (pre)fascio
\end{example}

\begin{remark}\label{mk_colimiti_filtrati_quoziente}\index{Colimite!filtrato!relazione di equivalenza}\index{Coequalizzatore!relazione di equivalenza}
 I colimiti filtrati sono importanti perché la relazione che genera il quoziente che definisce il coeq è già di equivalenza (grazie agli assiomi di cat filtrata, si apprezza particolarmente bene per gli ordini diretti...)
\end{remark}

\begin{example}\label{ex_completamenti_adici}\index{Completamento!adico}
 Completamenti adici
\end{example}

\begin{example}\label{ex_algebre_iniziali_coalgebre_terminali}\index{Algebra iniziale}\index{Coalgebra terminale}
 Algebre iniziali, coalgebre terminali
\end{example}

\begin{examples}\label{exs_algebre_iniziali_coalgebre_terminali}\index{Algebra iniziale!esempi}\index{Coalgebra terminale!esempi}
 Esempi di algebre iniziali e coalgebre terminali.
\end{examples}

\begin{theorem}\label{thm_adamek}\index{Teorema!di Adamek}\index{Adamek!teorema di}
 Teorema di Adamek
\end{theorem}
% \subsection{Limiti e colimiti in categorie concrete}
% \subsection{Colimiti in categorie algebriche}
\begin{esercizi}
 \item
 \item
 \item
 \item
 \item
\end{esercizi}
